% SPDX-License-Identifier: Apache-2.0
\section{An Array of Registers}
\label{sec:x-regs}

A unit over a count of ports often keeps something per port: a count,
a priority, whether a slot is full. \code{Regs<T, N>} is \code{N}
registers of one type as one field, and \code{self.counts[i]} is the
\code{i}-th, read with \code{get} and driven in \code{with!} as
\code{counts[i]: v} (issue 594). It is a type of its own rather than
an array of \code{Reg}, since a unit is \code{Default} and Rust has no
\code{Default} for an array of a generic length.

\code{Tally<N>} counts the hits on each of its input lines and puts
out the sum of the counts, going over the lines and their registers
in one loop. The netlist of \code{Tally<3>} has the registers
\code{counts\_0} to \code{counts\_2}, which the run prints, and it is
checked against the run under nvc and Verilator, the registers
included, as Section~\ref{sec:x-checked} describes. The index is a
number or the loop's variable: a register chosen by a signal would be
a multiplexer over all of them, and the lowering refuses it.

\begin{figure}[htbp]
\centering
\input{regs_timing.tex}
\caption{Three lines hit every cycle, every other and every third,
and the count each one keeps, a cycle behind its hit; the total is
the sum of the counts.}
\label{fig:x-regs}
\end{figure}

\lstinputlisting[caption={\code{ex\_regs.rs}. Run with \code{bazel run //lib/examples:ex\_regs}.},label={lst:x-regs}]{lib/examples/ex_regs.rs}

\lstinputlisting[style=ascii,caption={What \code{ex\_regs} printed when the build ran it: the totals, the registers, then the lowering.},label={lst:x-regs-out}]{out_regs.txt}
